% arquitectura_backend.tex — Capítulo 3: Arquitectura del Backend
\chapter{Arquitectura del Backend}

\section{Patrón Arquitectónico}

El backend de FINARA implementa una \textbf{Arquitectura por Capas} inspirada en el patrón MVC (Model-View-Controller), adaptada al contexto de una API REST donde no existe una capa de vistas. Las tres capas del sistema son:

\begin{enumerate}[leftmargin=1.5cm]
  \item \textbf{Capa de Enrutamiento (\texttt{routes/})}: recibe las solicitudes HTTP, aplica middleware de validación y delega al controlador correspondiente. No contiene lógica de negocio.
  \item \textbf{Capa de Controladores (\texttt{controllers/})}: orquesta la lógica de negocio. Consulta la base de datos, procesa los datos y construye la respuesta JSON.
  \item \textbf{Capa de Datos (\texttt{config/database.js})}: gestiona el pool de conexiones MySQL. Los controladores interactúan directamente con el pool mediante \textit{prepared statements}.
\end{enumerate}

Esta estructura fue elegida por su sencillez y adecuación al tamaño del proyecto. La adición de una capa de servicios o repositorios (\textit{Clean Architecture}) añadiría complejidad sin beneficio significativo para un equipo de dos personas \cite{martin_clean_arch}.

\section{Estructura de Carpetas Completa}

\begin{lstlisting}[style=bash, caption={Árbol de directorios del backend FINARA}, label={lst:tree}]
finara-backend/
  server.js                       # Entry point: carga .env, verifica BD, levanta Express
  .env                            # Variables de entorno (no versionado en Git)
  .env.example                    # Plantilla de variables de entorno
  package.json                    # Dependencias y scripts npm
  README.md                       # Guia de setup y ejecucion
  database/
    schema.sql                    # DDL completo: 10 tablas + 2 triggers
  ocr/
    LecturaTickets_v3.py          # Script OCR: pipeline adaptativo
    categorias.json               # 9 categorias con 200+ palabras clave
    tiendas.json                  # 35 cadenas comerciales con variantes OCR
  src/
    app.js                        # Configuracion Express: middleware, rutas, errores
    config/
      database.js                 # Pool MySQL con keepAlive y reconexion automatica
      upload.js                   # Configuracion Multer (tipo, tamanio, destino)
    controllers/
      authController.js           # Register, login, logout, verify, change-password
      categoryController.js       # CRUD categorias
      budgetController.js         # CRUD presupuestos + estado de uso
      incomeController.js         # CRUD ingresos + total mensual
      expenseController.js        # CRUD gastos + filtros + reversion en presupuesto
      ticketController.js         # OCR scan + CRUD tickets + productos
      debtController.js           # CRUD deudas + registro de pagos
      alertController.js          # CRUD alertas + config notificaciones
      dashboardController.js      # Resumen financiero + datos para graficas
      userController.js           # Perfil usuario + config inicial (onboarding)
    middleware/
      auth.js                     # Verificacion JWT en rutas protegidas
      validator.js                # Centralizacion de errores express-validator
      rateLimiter.js              # Rate limiting por contexto (auth, ocr, api)
      logger.js                   # Log coloreado de requests con tiempo de respuesta
      sanitizer.js                # Escapado HTML en todos los inputs del body
    routes/
      auth.js                     # /api/auth/*
      categories.js               # /api/categories/*
      budgets.js                  # /api/budgets/*
      incomes.js                  # /api/incomes/*
      expenses.js                 # /api/expenses/*
      tickets.js                  # /api/tickets/*
      debts.js                    # /api/debts/*
      alerts.js                   # /api/alerts/*
      dashboard.js                # /api/dashboard/*
      users.js                    # /api/users/*
    services/
      ocrService.js               # Invocacion del script Python + parse de JSON stdout
    utils/
      helpers.js                  # successResponse, errorResponse, suggestCategory, etc.
  tests/
    auth.test.js                  # 8 casos: register, login, verify
    debts.test.js                 # 9 casos: CRUD deudas, pago, liquidacion
    alerts.test.js                # 10 casos: alertas, config notificaciones
    helpers.test.js               # 12 casos: funciones utilitarias
    middleware.test.js            # 9 casos: sanitizer, auth JWT
  uploads/                        # Imagenes de tickets (creada automaticamente)
\end{lstlisting}

\section{Descripción de Módulos Principales}

\begin{longtable}{>{\ttfamily}p{4.5cm} p{9cm}}
  \caption{Responsabilidad de cada controlador}
  \label{tab:controllers} \\
  \toprule
  \textbf{Controlador} & \textbf{Responsabilidad} \\
  \midrule
  \endfirsthead
  \toprule
  \textbf{Controlador} & \textbf{Responsabilidad} \\
  \midrule
  \endhead
  authController      & Registro con hash bcrypt, login con comparación segura, generación/verificación JWT, cambio de contraseña \\
  \addlinespace
  categoryController  & CRUD de categorías predefinidas y personalizadas; bloqueo de eliminación si tienen gastos vinculados \\
  \addlinespace
  budgetController    & CRUD de presupuestos con cálculo de porcentaje utilizado; consulta por categoría \\
  \addlinespace
  incomeController    & CRUD de ingresos con filtro mensual y cálculo de total por período \\
  \addlinespace
  expenseController   & CRUD de gastos con reversión automática de \texttt{monto\_utilizado} al eliminar; filtros por mes, categoría y recientes \\
  \addlinespace
  ticketController    & Recibe imágenes multipart, invoca ocrService, crea Ticket + Productos\_Tickets + Gasto de forma atómica \\
  \addlinespace
  debtController      & CRUD de deudas con cálculo de \texttt{porcentaje\_pagado}; registro de pagos con avance de fecha y detección de liquidación \\
  \addlinespace
  alertController     & CRUD de alertas; marcar como leída (individual y masiva); CRUD de config de notificaciones \\
  \addlinespace
  dashboardController & Resumen financiero mensual (ingresos, gastos, balance, top categorías); datos para gráficas de gastos y presupuestos \\
  \addlinespace
  userController      & Perfil de usuario; actualización de flags de onboarding; configuración inicial (modo, distribución, ingreso base) \\
  \bottomrule
\end{longtable}

\section{Flujo de una Solicitud HTTP}

\begin{figure}[H]
  \centering
  \begin{tikzpicture}[
    node distance=0.7cm,
    box/.style={rectangle, rounded corners=3pt, draw=tableheadbg, fill=blue!5,
                minimum width=3.2cm, minimum height=0.8cm,
                font=\small, align=center},
    mid/.style={box, fill=orange!10},
    db/.style={cylinder, shape border rotate=90, draw=green!50!black, fill=green!5,
               minimum width=2cm, minimum height=0.7cm, align=center, font=\small},
    err/.style={box, fill=red!10},
    arr/.style={-Stealth, thick, tableheadbg}
  ]
    \node[box] (req)    {Request HTTP\\(Flutter / Postman)};
    \node[box, below=of req] (cors)   {CORS + Headers\\de Seguridad};
    \node[box, below=of cors] (logger) {Logger\\(método, ruta, user)};
    \node[box, below=of logger] (san)   {Sanitizer\\(escapado HTML)};
    \node[box, below=of san] (rate)   {Rate Limiter\\(auth / ocr / api)};
    \node[mid, below=of rate]  (route)  {Router\\(routes/*.js)};
    \node[mid, below=of route] (auth)   {Auth Middleware\\(JWT verify)};
    \node[mid, below=of auth]  (val)    {express-validator\\(validación body)};
    \node[mid, below=of val]   (ctrl)   {Controller\\(lógica de negocio)};
    \node[db,  below=of ctrl]  (mysql)  {MySQL Pool\\(prepared stmt)};
    \node[box, below=of mysql] (res)    {JSON Response\\(success/error)};

    \draw[arr] (req)    -- (cors);
    \draw[arr] (cors)   -- (logger);
    \draw[arr] (logger) -- (san);
    \draw[arr] (san)    -- (rate);
    \draw[arr] (rate)   -- (route);
    \draw[arr] (route)  -- (auth);
    \draw[arr] (auth)   -- (val);
    \draw[arr] (val)    -- (ctrl);
    \draw[arr] (ctrl)   -- (mysql);
    \draw[arr] (mysql)  -- (res);

    % Error paths
    \node[err, right=2cm of auth] (err401) {401\\Unauthorized};
    \node[err, right=2cm of val]  (err400) {400\\Bad Request};
    \node[err, right=2cm of rate] (err429) {429\\Too Many Req.};

    \draw[arr, red!70, dashed] (auth) -- (err401);
    \draw[arr, red!70, dashed] (val)  -- (err400);
    \draw[arr, red!70, dashed] (rate) -- (err429);
  \end{tikzpicture}
  \caption{Flujo de procesamiento de una solicitud HTTP en FINARA Backend}
  \label{fig:request-flow}
\end{figure}

\section{Formato Estándar de Respuestas}

Todos los endpoints de la API retornan JSON con la misma estructura, diferenciando entre éxito y error mediante el campo \texttt{success}.

\begin{lstlisting}[style=json, caption={Formatos de respuesta estándar de la API}, label={lst:response-format}]
// Respuesta exitosa
{
  "success": true,
  "data": { ... }   // objeto o arreglo segun el endpoint
}

// Respuesta de error
{
  "success": false,
  "error": "Mensaje descriptivo del error"
}

// Error de validacion (400)
{
  "success": false,
  "error": "El primer mensaje de error",
  "errors": [
    { "field": "email", "msg": "Email invalido" },
    { "field": "password", "msg": "Minimo 6 caracteres" }
  ]
}
\end{lstlisting}

\section{Flujo Especial: Escaneo OCR de Ticket}

El flujo del endpoint \texttt{POST /api/tickets/scan} involucra componentes adicionales:

\begin{lstlisting}[style=javascript, caption={Flujo completo de escaneo de ticket con OCR}, label={lst:ocr-flow}]
// 1. Multer recibe las imagenes (multipart/form-data)
//    Valida tipo (image/*) y tamanio (max 10MB)
//    Guarda en: uploads/<uuid>.jpg

// 2. ticketController.scan() invoca ocrService
const ocrResult = await processMultipleImages(uploadedPaths);

// 3. ocrService.js lanza el proceso Python
const proc = spawn('python3', ['ocr/LecturaTickets_v3.py', imagePath]);
// El script Python imprime JSON puro por stdout
// Los logs de debug van por stderr (invisibles para Node.js)

// 4. Python retorna JSON estructurado:
// { success, store_name, total, confidence, products: [...] }

// 5. ticketController inserta atomicamente:
//    INSERT INTO Tickets (...)
//    INSERT INTO Productos_Tickets (...)  -- uno por producto
//    INSERT INTO Gastos (...)             -- gasto automatico vinculado

// 6. Los triggers MySQL se disparan automaticamente:
//    - actualizar_presupuesto_gasto: actualiza monto_utilizado
//    - verificar_presupuesto_excedido: crea Alerta si >= 90%
\end{lstlisting}
